\chapter{Transposes and Permutations}

In Chapter~9, we watched Gaussian elimination run into an awkward obstacle: whenever
a zero landed in the pivot position on the diagonal, elimination could not divide by zero
and ground to a halt. The rescue was simple: swap two rows and keep going. We collected
those row exchanges into a single shuffle matrix $\mat{P}$, producing the complete factorization
$\mat{P}\mat{A} = \mat{L}\mat{U}$.

At the time, $\mat{P}$ seemed like mere emergency housekeeping. But row swaps and their
constant partner, the \emph{transpose}, are far from bookkeeping: they carry an enormous
amount of geometry. The first operation is the \emph{transpose}: flip a matrix across its
main diagonal so that rows become columns. The second is a \emph{permutation}: shuffle the
rows (or columns) of a matrix without changing any of the numbers, only their addresses.

The transpose is the operation that connects a matrix to the dot product. It is the reason
$\inner{\mat{A}\vec{x}}{\vec{y}} = \inner{\vec{x}}{\mat{A}\T\vec{y}}$, which is the single most
useful identity in all of applied linear algebra. It is the source of symmetric matrices and of
$\mat{A}\T\mat{A}$, the matrix behind least squares, covariance, and the singular value
decomposition. And permutations are our \emph{first} orthogonal matrices --- rigid rearrangements
of space that preserve every length and every angle. By the end of this chapter you will see the
transpose not as a typographical flip but as the geometric ``adjoint'': the matrix that does to the
output space exactly what $\mat{A}$ does to the input space, but in reverse.

Let us build all of this from the ground up, pictures first.

\section{The transpose: flipping across the diagonal}

\subsection{The definition, with a picture}

Start with the most concrete possible thing. Take a matrix and reflect it across the diagonal that
runs from the top-left corner down to the bottom-right. Whatever was in row $i$, column $j$ lands in
row $j$, column $i$.

\begin{definition}[Transpose]\label{ch10:def-transpose}
Let $\mat{A}$ be an $m\times n$ matrix with entries $a_{ij}$ (row $i$, column $j$). Its
\emph{transpose} $\mat{A}\T$ is the $n\times m$ matrix whose $(i,j)$ entry is $a_{ji}$:
\[
  (\mat{A}\T)_{ij} = a_{ji}.
\]
Rows of $\mat{A}$ become columns of $\mat{A}\T$, and columns of $\mat{A}$ become rows of $\mat{A}\T$.
\end{definition}

Here is a $2\times 3$ matrix and its $3\times 2$ transpose:
\[
  \mat{A} = \begin{bmatrix} 1 & 2 & 3 \\ 4 & 5 & 6 \end{bmatrix},
  \qquad
  \mat{A}\T = \begin{bmatrix} 1 & 4 \\ 2 & 5 \\ 3 & 6 \end{bmatrix}.
\]
Look at where each number went. The first \emph{row} of $\mat{A}$, namely $(1,2,3)$, became the first
\emph{column} of $\mat{A}\T$. The first column of $\mat{A}$, namely $\colvec{1\\4}$, became the first
row of $\mat{A}\T$. The diagonal entries that sit on the fold --- here only $a_{11}=1$ and $a_{22}=5$ ---
do not move. Everything else swaps across the fold.

\begin{intuitionbox}
Picture the matrix as a grid of tiles. Lay a mirror along the main diagonal (top-left to
bottom-right) and the transpose is what you see in the mirror. A tall thin matrix becomes a short
wide one. A row vector becomes a column vector. Nothing is added, scaled, or destroyed --- the same
numbers just trade places across the fold.
\end{intuitionbox}

\begin{figure}[H]
\centering
\begin{tikzpicture}[scale=0.85]
  % left grid A (2x3)
  \begin{scope}[shift={(0,0)}]
    \foreach \i in {0,1,2,3} \draw[gray!50] (\i,0) -- (\i,2);
    \foreach \j in {0,1,2} \draw[gray!50] (0,\j) -- (3,\j);
    % diagonal fold line drawn behind nodes
    \draw[trapred,very thick,dashed] (0,2) -- (2,0);
    \node[trapred,above right] at (0.1,2.05) {\small fold};
    \node[fill=white,inner sep=1.5pt] at (0.5,1.5) {$1$};
    \node[fill=white,inner sep=1.5pt] at (1.5,1.5) {$2$};
    \node[fill=white,inner sep=1.5pt] at (2.5,1.5) {$3$};
    \node[fill=white,inner sep=1.5pt] at (0.5,0.5) {$4$};
    \node[fill=white,inner sep=1.5pt] at (1.5,0.5) {$5$};
    \node[fill=white,inner sep=1.5pt] at (2.5,0.5) {$6$};
    \node[below] at (1.5,-0.15) {$\mat{A}$ \footnotesize ($2\times 3$)};
  \end{scope}
  % arrow
  \draw[-{Latex[length=3mm]},thick,sourceblue] (3.6,1) -- (5.0,1)
    node[midway,above]{\small transpose};
  % right grid A^T (3x2)
  \begin{scope}[shift={(5.6,-0.5)}]
    \foreach \i in {0,1,2} \draw[gray!50] (\i,0) -- (\i,3);
    \foreach \j in {0,1,2,3} \draw[gray!50] (0,\j) -- (2,\j);
    % diagonal fold line drawn behind nodes
    \draw[trapred,very thick,dashed] (0,3) -- (2,1);
    \node[trapred,above right] at (0.1,3.05) {\small fold};
    \node[fill=white,inner sep=1.5pt] at (0.5,2.5) {$1$};
    \node[fill=white,inner sep=1.5pt] at (1.5,2.5) {$4$};
    \node[fill=white,inner sep=1.5pt] at (0.5,1.5) {$2$};
    \node[fill=white,inner sep=1.5pt] at (1.5,1.5) {$5$};
    \node[fill=white,inner sep=1.5pt] at (0.5,0.5) {$3$};
    \node[fill=white,inner sep=1.5pt] at (1.5,0.5) {$6$};
    \node[below] at (1,-0.15) {$\mat{A}\T$ \footnotesize ($3\times 2$)};
  \end{scope}
\end{tikzpicture}
\caption{The transpose reflects the grid across the main diagonal. Row $(1,2,3)$ of $\mat{A}$
becomes the first column of $\mat{A}\T$. The shape flips from $2\times 3$ to $3\times 2$.}
\label{ch10:fig-flip}
\end{figure}

A few immediate sanity checks. If $\mat{A}$ is $m\times n$ then $\mat{A}\T$ is $n\times m$ --- the
shape flips. Transposing twice brings you home: $(\mat{A}\T)\T = \mat{A}$, because reflecting across
the diagonal and then reflecting again undoes itself. And a single column vector $\vec{x}$ (which we
think of as $n\times 1$) transposes into a single \emph{row}, $\vec{x}\T$, of shape $1\times n$. That
last fact is the key, because a row times a column is exactly a dot product.

\begin{focusbox}
Hold onto three reflexes:
\begin{itemize}
  \item $(\mat{A}\T)_{ij} = a_{ji}$: swap the indices.
  \item Shape flips: $m\times n \to n\times m$.
  \item $(\mat{A}\T)\T = \mat{A}$: transpose is its own undo.
\end{itemize}
\end{focusbox}

\begin{recoverybox}
\textbf{Pause and think:} If a matrix $\mat{M}$ has shape $4\times 7$, what is the shape of $(\mat{M}\T)\T$?
Which entries of $\mat{M}$ never change position when you transpose?
\emph{Answer:} Transposing once produces $7\times 4$; transposing again returns the original $4\times 7$.
Only the diagonal entries $m_{11}, m_{22}, m_{33}, m_{44}$ stay glued in place, because swapping row index $i$
with column index $i$ leaves the address unchanged.
\end{recoverybox}

\subsection{The transpose turns a column into a dot-product machine}

Why should anyone care about flipping a grid? Here is the first real reason. Recall the dot product
of two column vectors $\vec{u},\vec{v}\in\R^n$:
\[
  \inner{\vec{u}}{\vec{v}} \;=\; u_1 v_1 + u_2 v_2 + \cdots + u_n v_n.
\]
Now write $\vec{u}$ as a column and transpose it into a row. Matrix-multiply that row by the column
$\vec{v}$:
\[
  \vec{u}\T \vec{v}
  = \begin{bmatrix} u_1 & u_2 & \cdots & u_n \end{bmatrix}
    \begin{bmatrix} v_1 \\ v_2 \\ \vdots \\ v_n \end{bmatrix}
  = u_1v_1 + u_2v_2 + \cdots + u_nv_n
  = \inner{\vec{u}}{\vec{v}}.
\]
So the dot product \emph{is} a transpose-times-a-vector. The little operation ``$\vec{u}\T\vec{v}$''
is matrix multiplication of a $1\times n$ row by an $n\times 1$ column, producing a $1\times 1$
number. This is more than a notational convenience; it lets every dot-product fact be
written as a matrix product and back again. We will lean on it constantly.

\begin{example}[A dot product as $\vec{u}\T\vec{v}$]\label{ch10:ex-dot}
Let $\vec{u}=\colvec{2\\-1\\3}$ and $\vec{v}=\colvec{1\\4\\2}$. Then
\[
  \vec{u}\T\vec{v}
  = \begin{bmatrix} 2 & -1 & 3 \end{bmatrix}\colvec{1\\4\\2}
  = (2)(1)+(-1)(4)+(3)(2) = 2-4+6 = 4.
\]
And $\norm{\vec{u}}^2 = \vec{u}\T\vec{u} = 4+1+9 = 14$. A vector dotted with itself, written
$\vec{u}\T\vec{u}$, is the squared length. Remember that --- $\mat{A}\T\mat{A}$ later is just this
idea applied to every pair of columns at once.

Notice the order of multiplication carefully. Row times column, $(1\times n)(n\times 1)$, produces
a $1\times 1$ scalar: the inner product (dot product). If you reverse the order to column times row,
$(n\times 1)(1\times n)$, you obtain an $n\times n$ grid of numbers: the \emph{outer product} $\vec{u}\vec{v}\T$.
Every row of an outer product is a scaled copy of $\vec{v}\T$, and every column is a scaled copy of $\vec{u}$.
Keep row-times-column (number) sharply distinct from column-times-row (rank-one matrix).
\end{example}

\section{The algebra of transposes}

Before we go further geometrically we need the rulebook. There are four rules; three are easy and
one is the famous ``socks and shoes'' reversal. We will derive each, because a formula you can
derive in three lines is a formula you will never misremember.

\begin{theorem}[Rules of the transpose]\label{ch10:thm-rules}
For matrices of compatible sizes and a scalar $c$:
\begin{enumerate}
  \item $(\mat{A}+\mat{B})\T = \mat{A}\T + \mat{B}\T$.
  \item $(c\,\mat{A})\T = c\,\mat{A}\T$.
  \item $(\mat{A}\T)\T = \mat{A}$.
  \item $(\mat{A}\mat{B})\T = \mat{B}\T\mat{A}\T$ \quad(the order reverses!).
\end{enumerate}
\end{theorem}

\begin{proof}
Rules 1--3 are entry-chasing. For (1), the $(i,j)$ entry of $(\mat{A}+\mat{B})\T$ is by definition the
$(j,i)$ entry of $\mat{A}+\mat{B}$, which is $a_{ji}+b_{ji}$, which is the $(i,j)$ entry of
$\mat{A}\T+\mat{B}\T$. Rule (2) is the same with a scalar pulled along. Rule (3) we already saw:
reflecting twice is the identity.

Rule (4) is the one worth doing carefully. Let $\mat{A}$ be $m\times n$ and $\mat{B}$ be $n\times p$,
so $\mat{A}\mat{B}$ is $m\times p$ and $(\mat{A}\mat{B})\T$ is $p\times m$. The $(i,j)$ entry of
$(\mat{A}\mat{B})\T$ is the $(j,i)$ entry of $\mat{A}\mat{B}$:
\[
  \big((\mat{A}\mat{B})\T\big)_{ij}
  = (\mat{A}\mat{B})_{ji}
  = \sum_{k=1}^{n} a_{jk}\, b_{ki}.
\]
Now compute the $(i,j)$ entry of $\mat{B}\T\mat{A}\T$. The $(i,k)$ entry of $\mat{B}\T$ is $b_{ki}$,
and the $(k,j)$ entry of $\mat{A}\T$ is $a_{jk}$, so
\[
  (\mat{B}\T\mat{A}\T)_{ij}
  = \sum_{k=1}^{n} (\mat{B}\T)_{ik}(\mat{A}\T)_{kj}
  = \sum_{k=1}^{n} b_{ki}\, a_{jk}.
\]
The two sums are identical (multiplication of scalars commutes), so the matrices are equal.
\end{proof}

\begin{trapbox}
The order really does reverse. $(\mat{A}\mat{B})\T = \mat{B}\T\mat{A}\T$, \emph{not}
$\mat{A}\T\mat{B}\T$. A quick way to see why it \emph{has} to reverse: if $\mat{A}$ is $2\times 3$ and
$\mat{B}$ is $3\times 4$, then $\mat{A}\T\mat{B}\T$ tries to multiply a $3\times 2$ by a $4\times 3$ ---
the inner dimensions $2$ and $4$ do not match, so it is not even defined. Only $\mat{B}\T\mat{A}\T$
($4\times 3$ times $3\times 2$) makes sense.
\end{trapbox}

A standard mnemonic illustrates the reversal: to take off your socks and shoes you reverse the order in which you put
them on: shoes first, then socks. Transposing a product is the same --- whatever went on last comes
off first. Combined with the inverse rule $(\mat{A}\mat{B})\inv = \mat{B}\inv\mat{A}\inv$ from the
previous chapter, you get a clean fact: transpose and inverse commute with each other,
$(\mat{A}\T)\inv = (\mat{A}\inv)\T$, a matrix so common it gets its own name, $\mat{A}^{-\mathsf{T}}$.

\begin{labbox}
Verify the reversal rule directly. Take
\[
  \mat{A}=\begin{bmatrix}1&2\\0&1\end{bmatrix},\qquad
  \mat{B}=\begin{bmatrix}3&0\\4&5\end{bmatrix}.
\]
Compute $\mat{A}\mat{B}$, then transpose it. Separately compute $\mat{B}\T\mat{A}\T$. They must
match. Then compute the wrong thing, $\mat{A}\T\mat{B}\T$, and watch it disagree. (You should find
$(\mat{A}\mat{B})\T = \begin{bsmallmatrix}11&4\\10&5\end{bsmallmatrix}$.)
\end{labbox}

Why does this reversal happen geometrically? Think of $\mat{A}\mat{B}\vec{x}$ as a two-stage assembly
line: first $\mat{B}$ transforms $\vec{x}$, and then $\mat{A}$ transforms $\mat{B}\vec{x}$.
When we want to pull a measurement backward from the final output space back to the starting space
(as we are about to see), we must pull it backward through $\mat{A}$ first, and only then backward
through $\mat{B}$. The last operation to act in the forward direction is the first that must be
undone in the backward direction.

\begin{recoverybox}
\textbf{Pause and think:} What is the transpose of a three-matrix product $(\mat{A}\mat{B}\mat{C})\T$?
\emph{Answer:} Group the first two factors: $((\mat{A}\mat{B})\mat{C})\T = \mat{C}\T(\mat{A}\mat{B})\T = \mat{C}\T\mat{B}\T\mat{A}\T$.
The reversal rule applies to products of any length: reverse the factor order completely and transpose each matrix.
\end{recoverybox}

\section{The transpose is the adjoint}

Now comes the identity that explains \emph{why} the transpose is the right notion and not just a
formatting trick.

\begin{theorem}[The transpose moves across the inner product]\label{ch10:thm-adjoint}
Let $\mat{A}$ be an $m\times n$ real matrix, $\vec{x}\in\R^n$, and $\vec{y}\in\R^m$. Then
\[
  \inner{\mat{A}\vec{x}}{\vec{y}} \;=\; \inner{\vec{x}}{\mat{A}\T\vec{y}}.
\]
\end{theorem}

\begin{proof}
Write both sides as row-times-column products and use the reversal rule. The left side is
\[
  \inner{\mat{A}\vec{x}}{\vec{y}}
  = (\mat{A}\vec{x})\T\vec{y}
  = \vec{x}\T\mat{A}\T\vec{y},
\]
using rule (4): $(\mat{A}\vec{x})\T = \vec{x}\T\mat{A}\T$. The right side is
\[
  \inner{\vec{x}}{\mat{A}\T\vec{y}}
  = \vec{x}\T(\mat{A}\T\vec{y})
  = \vec{x}\T\mat{A}\T\vec{y}.
\]
Both equal the same $1\times 1$ quantity $\vec{x}\T\mat{A}\T\vec{y}$. Done.
\end{proof}

The proof is two lines, but its meaning is large. The matrix $\mat{A}$ takes vectors from the input
space $\R^n$ over to the output space $\R^m$. The
transpose $\mat{A}\T$ runs the other way, from $\R^m$ back to $\R^n$. The theorem says: it does not
matter whether you push $\vec{x}$ forward through $\mat{A}$ and then measure it against $\vec{y}$ in
the output space, or pull $\vec{y}$ backward through $\mat{A}\T$ and then measure it against $\vec{x}$
in the input space. You get the same number. The transpose is the unique matrix that lets you
``hand the matrix across'' the inner product.

\begin{figure}[H]
\centering
\begin{tikzpicture}[scale=1.0]
  % input space
  \node[draw,sourceblue,thick,rounded corners,minimum width=3.2cm,minimum height=3.4cm] (in)
    at (0,0) {};
  \node[sourceblue,below] at (0,-1.9) {input space $\R^n$};
  \node[sourceblue] (x) at (0,0.7) {$\vec{x}$};
  \node[sourceblue] (ATy) at (0,-0.7) {$\mat{A}\T\vec{y}$};
  % output space
  \node[draw,targetgreen,thick,rounded corners,minimum width=3.2cm,minimum height=3.4cm] (out)
    at (5,0) {};
  \node[targetgreen,below] at (5,-1.9) {output space $\R^m$};
  \node[targetgreen] (Ax) at (5,0.7) {$\mat{A}\vec{x}$};
  \node[targetgreen] (y) at (5,-0.7) {$\vec{y}$};
  % forward A
  \draw[-{Latex[length=2.5mm]},thick,sourceblue] (x) to[bend left=18] node[above]{$\mat{A}$} (Ax);
  % backward A^T
  \draw[-{Latex[length=2.5mm]},thick,targetgreen] (y) to[bend left=18] node[below]{$\mat{A}\T$} (ATy);
  % the two inner products
  \draw[densely dotted,quotientorange,thick] (Ax) -- (y)
    node[midway,right]{$\inner{\mat{A}\vec{x}}{\vec{y}}$};
  \draw[densely dotted,quotientorange,thick] (x) -- (ATy)
    node[midway,left]{$\inner{\vec{x}}{\mat{A}\T\vec{y}}$};
\end{tikzpicture}
\caption{The adjoint relation. Push $\vec{x}$ forward by $\mat{A}$ and pair with $\vec{y}$ on the
right; or pull $\vec{y}$ back by $\mat{A}\T$ and pair with $\vec{x}$ on the left. The two dotted
pairings give the same number.}
\label{ch10:fig-adjoint}
\end{figure}

\begin{aestheticbox}
This is the reason mathematicians often call $\mat{A}\T$ the \emph{adjoint} of $\mat{A}$. In more
advanced settings (infinite-dimensional spaces, differential operators, quantum mechanics) the
matrix entries vanish entirely and the \emph{only} thing that survives is exactly this relation:
the adjoint $\mat{A}^{*}$ of an operator $\mat{A}$ is \emph{defined} as the operator satisfying
$\inner{\mat{A}\vec{x}}{\vec{y}}=\inner{\vec{x}}{\mat{A}^{*}\vec{y}}$ for all $\vec{x},\vec{y}$. The
matrix transpose is the finite-dimensional foundation of this general concept in analysis. When you meet
$\frac{d}{dx}$ and its adjoint $-\frac{d}{dx}$ later, you will recognize this same picture.
\end{aestheticbox}

\begin{strangbox}
\textbf{Before Strang:} when he writes $\mat{A}\T$, do not read it as ``the flipped matrix.'' Read
it as ``the matrix that takes outputs back to inputs while preserving dot products.'' That reading
makes the four fundamental subspaces (coming a few chapters from now) feel inevitable rather than
mysterious: the column space of $\mat{A}\T$ is the row space of $\mat{A}$, and orthogonality between
the subspaces is just Theorem~\ref{ch10:thm-adjoint} with one side set to zero. \textbf{After
Strang:} this identity is the engine of backpropagation, of the normal equations, and of the
adjoint method in optimization and PDEs.
\end{strangbox}

\begin{applicationbox}
\textbf{Backpropagation is a transpose.} A layer of a neural network computes
$\vec{z}=\mat{W}\vec{x}$ (followed by a nonlinearity we will ignore here). Training needs the
gradient of a scalar loss $L$ with respect to the input, given the gradient with respect to the
output. The chain rule says: if $\vec{g}=\partial L/\partial \vec{z}$ is the ``upstream gradient'' in
the output space, then the gradient in the input space is
\[
  \frac{\partial L}{\partial \vec{x}} = \mat{W}\T \vec{g}.
\]
The forward pass multiplies by $\mat{W}$; the backward pass multiplies by $\mat{W}\T$. That is
exactly Theorem~\ref{ch10:thm-adjoint}: gradients are covectors that live in the output space, and
the transpose is what pulls them back to the input space. Every deep-learning framework implements
``the backward of a matmul is a matmul by the transpose'' --- the whole training loop is the adjoint
relation applied layer by layer.
\end{applicationbox}

\begin{recoverybox}
\textbf{Pause and think:} If $\mat{A}$ is a $3\times 2$ matrix, what spaces do $\vec{x}$ and $\vec{y}$
live in for the adjoint identity $\inner{\mat{A}\vec{x}}{\vec{y}} = \inner{\vec{x}}{\mat{A}\T\vec{y}}$?
In which space is each dot product evaluated?
\emph{Answer:} For $\mat{A}\vec{x}$ to make sense, $\vec{x}\in\R^2$, so the output $\mat{A}\vec{x}\in\R^3$.
Thus $\vec{y}\in\R^3$. The left-hand dot product $\inner{\mat{A}\vec{x}}{\vec{y}}$ takes place in the
$3$-dimensional output space $\R^3$. On the right-hand side, $\mat{A}\T$ has shape $2\times 3$, so
$\mat{A}\T\vec{y}\in\R^2$; the dot product $\inner{\vec{x}}{\mat{A}\T\vec{y}}$ takes place in the
$2$-dimensional input space $\R^2$. Two dot products in two different dimensions, yet the resulting
number is identical.
\end{recoverybox}

\section{Symmetric matrices}

Now that we understand how the transpose moves operators across inner products, a natural question
arises: what happens when transposing a matrix does not change it at all?

\subsection{When a matrix equals its own mirror image}

A matrix that is unchanged by transposing is the most important special case in the entire subject.

\begin{definition}[Symmetric matrix]\label{ch10:def-symmetric}
A square matrix $\mat{A}$ is \emph{symmetric} if $\mat{A}\T=\mat{A}$, that is, $a_{ij}=a_{ji}$ for all
$i,j$. Geometrically, it is unchanged when reflected across its main diagonal.
\end{definition}

For example,
\[
  \mat{S}=\begin{bmatrix} 2 & 7 & -1 \\ 7 & 0 & 3 \\ -1 & 3 & 5 \end{bmatrix}
\]
is symmetric: the entry above the diagonal at position $(1,2)$ is $7$, and so is the entry below at
$(2,1)$. The whole matrix is determined by its diagonal and the triangle above it; the lower triangle
is just a copy folded over. An $n\times n$ symmetric matrix has only $n + \binom{n}{2}=\binom{n+1}{2}$
free numbers instead of $n^2$.

\begin{intuitionbox}
A symmetric matrix is one where ``the influence of $i$ on $j$ equals the influence of $j$ on $i$.''
Think of a spring connecting two masses: the force the first exerts on the second equals the force
the second exerts on the first. Or a road network where the distance from town $i$ to town $j$ is
the same as from $j$ to $i$. Symmetry encodes a two-way, reciprocal relationship --- and it is no
accident that such matrices have the nicest possible eigenvectors (real, orthogonal), a fact we will
prove in the spectral-theorem chapter.
\end{intuitionbox}

\subsection{\texorpdfstring{$\mat{A}\T\mat{A}$}{ATA} is always symmetric --- and it is everywhere}

Here is a manufacturing process for symmetric matrices that works on \emph{any} matrix, square or
not. Take any $m\times n$ matrix $\mat{A}$ and form $\mat{A}\T\mat{A}$.

\begin{proposition}[The fundamental symmetric matrices]\label{ch10:prop-ATA}
For any $m\times n$ matrix $\mat{A}$, both $\mat{A}\T\mat{A}$ ($n\times n$) and $\mat{A}\mat{A}\T$
($m\times m$) are symmetric.
\end{proposition}

\begin{proof}
Apply the transpose rules. Using $(\mat{X}\mat{Y})\T=\mat{Y}\T\mat{X}\T$ and $(\mat{X}\T)\T=\mat{X}$,
\[
  (\mat{A}\T\mat{A})\T
  = \mat{A}\T(\mat{A}\T)\T
  = \mat{A}\T\mat{A}.
\]
Since $\mat{A}\T\mat{A}$ equals its own transpose, it is symmetric. The same computation with the
roles swapped gives $(\mat{A}\mat{A}\T)\T=\mat{A}\mat{A}\T$.
\end{proof}

Work one out and see what the entries mean.

\begin{example}[$\mat{A}\T\mat{A}$ collects all the dot products of the columns]\label{ch10:ex-ATA}
Let
\[
  \mat{A}=\begin{bmatrix} 1 & 0 \\ 1 & 1 \\ 0 & 2 \end{bmatrix}
  \quad\text{with columns}\quad
  \vec{a}_1=\colvec{1\\1\\0},\ \vec{a}_2=\colvec{0\\1\\2}.
\]
Then
\[
  \mat{A}\T\mat{A}
  = \begin{bmatrix} 1 & 1 & 0 \\ 0 & 1 & 2 \end{bmatrix}
    \begin{bmatrix} 1 & 0 \\ 1 & 1 \\ 0 & 2 \end{bmatrix}
  = \begin{bmatrix} 2 & 1 \\ 1 & 5 \end{bmatrix}.
\]
Read the entries: the $(1,1)$ entry $2=\vec{a}_1\T\vec{a}_1=\norm{\vec{a}_1}^2$; the $(2,2)$ entry
$5=\norm{\vec{a}_2}^2$; the off-diagonal $1=\vec{a}_1\T\vec{a}_2=\vec{a}_2\T\vec{a}_1$. In general the
$(i,j)$ entry of $\mat{A}\T\mat{A}$ is the dot product $\vec{a}_i\T\vec{a}_j$ of column $i$ with column
$j$. The matrix is symmetric precisely because $\vec{a}_i\T\vec{a}_j=\vec{a}_j\T\vec{a}_i$. This grid
of all pairwise dot products is so useful it has a name: the \emph{Gram matrix} of the columns.
\end{example}

\begin{focusbox}
$\mat{A}\T\mat{A}$ is the table of all dot products between the columns of $\mat{A}$. Diagonal entries
are squared lengths; off-diagonal entries measure how aligned two columns are (zero means
perpendicular). Symmetry is automatic because dot products do not care about order
(Figure~\ref{ch10:fig-gram}).
\end{focusbox}

\begin{figure}[H]
\centering
\begin{tikzpicture}[scale=1.0,every node/.style={font=\small}]
  % fold line across diagonal in 3 segments (outside and between diagonal entries)
  \draw[trapred,densely dashed,very thick] (-1.3,1.3) -- (-0.85,0.85);
  \draw[trapred,densely dashed,very thick] (-0.25,0.25) -- (0.25,-0.25);
  \draw[trapred,densely dashed,very thick] (0.85,-0.85) -- (1.3,-1.3);
  \node[trapred, above right] at (-1.1,1.1) {\footnotesize fold};

  % each column dotted with each column -> grid
  \matrix (m) [matrix of math nodes, nodes={minimum size=9mm}, column sep=4pt, row sep=4pt,
               left delimiter={[}, right delimiter={]}] at (0,0)
  {
    \vec{a}_1\T\vec{a}_1 & \vec{a}_1\T\vec{a}_2 \\
    \vec{a}_2\T\vec{a}_1 & \vec{a}_2\T\vec{a}_2 \\
  };
  \node[below=5mm of m] {grid of all pairwise dot products $=\mat{A}\T\mat{A}$};
  % equals the numbers
  \node[right=8mm of m] (eq) {$=$};
  \matrix (n) [matrix of math nodes, nodes={minimum size=9mm}, column sep=4pt, row sep=4pt,
               left delimiter={[}, right delimiter={]}, right=2mm of eq]
  {
    2 & 1 \\
    1 & 5 \\
  };
\end{tikzpicture}
\caption{The Gram matrix from Example~\ref{ch10:ex-ATA}. Entry $(i,j)$ is $\vec{a}_i\T\vec{a}_j$, the
dot product of column $i$ with column $j$. Reflecting across the dashed diagonal leaves it unchanged
because $\vec{a}_i\T\vec{a}_j=\vec{a}_j\T\vec{a}_i$: the matrix is symmetric by construction.}
\label{ch10:fig-gram}
\end{figure}

\begin{applicationbox}
\textbf{Covariance matrices and PCA.} Suppose you have a data set: $m$ samples, each with $n$
measured features, stacked as the rows of a matrix $\mat{X}$ (so $\mat{X}$ is $m\times n$). After
subtracting the mean of each column, the matrix
\[
  \mat{C} = \tfrac{1}{m-1}\,\mat{X}\T\mat{X}
\]
is the \emph{covariance matrix} of the features. Its $(i,j)$ entry is the covariance between feature
$i$ and feature $j$; its diagonal holds the variances. It is symmetric for exactly the reason in
Example~\ref{ch10:ex-ATA}, and that symmetry guarantees real eigenvalues and orthogonal eigenvectors
--- the principal components. Principal Component Analysis, the core tool for dimensionality
reduction, is exactly the spectral analysis of a symmetric $\mat{X}\T\mat{X}$. The transpose is
what turns raw data into reciprocal, analyzable structure.
\end{applicationbox}

\subsection{Splitting any square matrix into symmetric plus antisymmetric}

A nice structural fact, and a good exercise in the rules. A matrix with $\mat{A}\T=-\mat{A}$ is called
\emph{antisymmetric} (or skew-symmetric); its diagonal must be zero and entries flip sign across the
fold. Every square matrix splits uniquely into a symmetric part and an antisymmetric part:
\[
  \mat{A}
  = \underbrace{\tfrac12(\mat{A}+\mat{A}\T)}_{\text{symmetric}}
  + \underbrace{\tfrac12(\mat{A}-\mat{A}\T)}_{\text{antisymmetric}}.
\]
Check the first piece is symmetric: $\big(\tfrac12(\mat{A}+\mat{A}\T)\big)\T
=\tfrac12(\mat{A}\T+\mat{A})$, the same thing. Check the second is antisymmetric:
$\big(\tfrac12(\mat{A}-\mat{A}\T)\big)\T=\tfrac12(\mat{A}\T-\mat{A})$, which is the negative of the
original. Add the two pieces and the $\mat{A}\T$ terms cancel, returning $\mat{A}$.

\begin{connectionbox}
This decomposition is the matrix mirror of even-plus-odd in calculus. Any function splits as
$f(x)=\tfrac12\big(f(x)+f(-x)\big)+\tfrac12\big(f(x)-f(-x)\big)$, an even part plus an odd part,
exactly as a matrix splits into a transpose-even (symmetric) part plus a transpose-odd
(antisymmetric) part. In both cases an \emph{involution} --- an operation that serves as its own undo
when applied twice, like negating an input ($x \mapsto -x$) or reflecting across a mirror --- carves
the space into a $+1$ piece and a $-1$ piece. The antisymmetric $3\times 3$
matrices, by the way, are exactly the cross-product operators $\vec{v}\mapsto \vec{w}\times\vec{v}$,
which is why rotations in $\R^3$ come from antisymmetric generators. Symmetry and antisymmetry under
the transpose are a recurring theme you will keep meeting.
\end{connectionbox}

\section{Permutation matrices: shuffling without distorting}

\subsection{Reordering the rows of the identity}

We now switch to the second star of the chapter. A permutation matrix is what you get by taking the
identity matrix and rearranging its rows. Nothing more. Yet these matrices are the cleanest possible
example of an \emph{orthogonal} matrix, our first rigid motion of space.

\begin{definition}[Permutation matrix]\label{ch10:def-perm}
A \emph{permutation matrix} $\mat{P}$ is a square matrix obtained by permuting the rows (equivalently,
the columns) of the identity $\mat{I}$. Equivalently, $\mat{P}$ has exactly one $1$ in each row and
each column and zeros elsewhere.
\end{definition}

The smallest nontrivial example swaps two coordinates. In $\R^2$,
\[
  \mat{P}=\begin{bmatrix} 0 & 1 \\ 1 & 0 \end{bmatrix},
  \qquad
  \mat{P}\colvec{x\\y}=\colvec{y\\x}.
\]
It takes the point $(x,y)$ to $(y,x)$ --- a reflection across the diagonal line $y=x$. In $\R^3$ a
permutation can do something more interesting, a \emph{cycle}:
\[
  \mat{P}=\begin{bmatrix} 0 & 0 & 1 \\ 1 & 0 & 0 \\ 0 & 1 & 0 \end{bmatrix},
  \qquad
  \mat{P}\colvec{x\\y\\z}=\colvec{z\\x\\y}.
\]
The old third coordinate becomes the new first, the old first becomes the new second, and the old
second becomes the new third: a $3$-cycle. Notice that whatever $\mat{P}$ does, it never adds,
scales, or blends coordinates. It only relocates them. A permutation does only one thing:
it rigidly relabels the coordinate axes.

Because $\mat{P}\vec{e}_j$ is the $j$th column of $\mat{P}$, the action on the standard basis is the
fastest way to read a permutation. Here the columns of $\mat{P}$ are $\vec{e}_2,\vec{e}_3,\vec{e}_1$,
so $\vec{e}_1\mapsto\vec{e}_2\mapsto\vec{e}_3\mapsto\vec{e}_1$: the three axes chase each other around
a single loop (Figure~\ref{ch10:fig-cycle}).

\begin{figure}[H]
\centering
\begin{tikzpicture}[scale=1.0,every node/.style={font=\small}]
  \node[sourceblue,circle,draw,thick,minimum size=8mm] (e1) at (90:1.7) {$\vec{e}_1$};
  \node[sourceblue,circle,draw,thick,minimum size=8mm] (e2) at (-30:1.7) {$\vec{e}_2$};
  \node[sourceblue,circle,draw,thick,minimum size=8mm] (e3) at (210:1.7) {$\vec{e}_3$};
  \draw[-{Latex[length=2.5mm]},targetgreen,very thick] (e1) to[bend left=22] (e2);
  \draw[-{Latex[length=2.5mm]},targetgreen,very thick] (e2) to[bend left=22] (e3);
  \draw[-{Latex[length=2.5mm]},targetgreen,very thick] (e3) to[bend left=22] (e1);
  \node[targetgreen] at (0,0) {$\mat{P}$};
\end{tikzpicture}
\caption{The $3$-cycle reads off the columns of $\mat{P}$: column $j$ is $\mat{P}\vec{e}_j$. Here
$\vec{e}_1\mapsto\vec{e}_2\mapsto\vec{e}_3\mapsto\vec{e}_1$, one closed loop. The transpose
$\mat{P}\T$ runs every arrow backward, which is why it undoes the shuffle.}
\label{ch10:fig-cycle}
\end{figure}

\begin{intuitionbox}
A permutation matrix is a perfect shuffle of the coordinate axes. The standard basis vectors
$\vec{e}_1,\dots,\vec{e}_n$ get permuted among themselves and sent back to standard basis vectors ---
none is bent or stretched. The unit cube goes to the unit cube; only the labels on the corners move.
This is why permutations preserve every length and every angle.
\end{intuitionbox}

\begin{figure}[H]
\centering
\begin{tikzpicture}[scale=1.1]
  % left: original square with labeled corners
  \begin{scope}
    \draw[thin,gray!40] (-0.2,0)--(2.2,0);
    \draw[thin,gray!40] (0,-0.2)--(0,2.2);
    \fill[sourceblue!12] (0,0)--(2,0)--(2,2)--(0,2)--cycle;
    \draw[sourceblue,thick] (0,0)--(2,0)--(2,2)--(0,2)--cycle;
    \draw[-{Latex[length=2mm]},vecblue,thick] (0,0)--(2,0) node[below right]{$\vec{e}_1$};
    \draw[-{Latex[length=2mm]},vecred,thick] (0,0)--(0,2) node[above left]{$\vec{e}_2$};
    \node[below] at (1,-0.5) {before};
  \end{scope}
  % arrow
  \draw[-{Latex[length=3mm]},thick] (2.7,1)--(4.0,1)
    node[midway,above]{\small $\mat{P}=\left[\begin{smallmatrix}0&1\\1&0\end{smallmatrix}\right]$};
  % right: swapped axes
  \begin{scope}[shift={(4.6,0)}]
    \draw[thin,gray!40] (-0.2,0)--(2.2,0);
    \draw[thin,gray!40] (0,-0.2)--(0,2.2);
    \fill[sourceblue!12] (0,0)--(2,0)--(2,2)--(0,2)--cycle;
    \draw[sourceblue,thick] (0,0)--(2,0)--(2,2)--(0,2)--cycle;
    \draw[-{Latex[length=2mm]},vecred,thick] (0,0)--(2,0) node[below right]{$\vec{e}_2$};
    \draw[-{Latex[length=2mm]},vecblue,thick] (0,0)--(0,2) node[above left]{$\vec{e}_1$};
    \draw[densely dashed,trapred] (-0.2,-0.2)--(2.3,2.3)
      node[above right=-1pt]{\small $y=x$};
    \node[below] at (1,-0.5) {after (axes swapped)};
  \end{scope}
\end{tikzpicture}
\caption{The swap $\mat{P}=\left[\begin{smallmatrix}0&1\\1&0\end{smallmatrix}\right]$ exchanges the
two axes. It is a reflection across the line $y=x$: the square is unchanged as a set, but $\vec{e}_1$
and $\vec{e}_2$ trade places.}
\label{ch10:fig-swap}
\end{figure}

\subsection{Multiplying by $\mat{P}$ on the left vs. on the right}

Remember the handedness of this action: multiplying \emph{on the left} by $\mat{P}$ permutes the \emph{rows} of
whatever it hits; multiplying \emph{on the right} permutes the \emph{columns}. Watch it with the swap
matrix:
\[
  \begin{bmatrix}0&1\\1&0\end{bmatrix}
  \begin{bmatrix}a&b\\c&d\end{bmatrix}
  = \begin{bmatrix}c&d\\a&b\end{bmatrix}
  \quad(\text{rows swapped}),
  \qquad
  \begin{bmatrix}a&b\\c&d\end{bmatrix}
  \begin{bmatrix}0&1\\1&0\end{bmatrix}
  = \begin{bmatrix}b&a\\d&c\end{bmatrix}
  \quad(\text{columns swapped}).
\]
This is exactly the row-exchange operation you met during elimination: when a pivot is zero you swap
rows, and that swap is recorded as a left multiplication by a permutation matrix. (In $\mat{P}\mat{A}=\mat{L}\mat{U}$
the $\mat{P}$ is precisely the bookkeeping of which rows got exchanged.)

\subsection{How many are there, and how do they multiply?}

There are exactly $n!$ permutation matrices of size $n\times n$, one for each way to reorder $n$
items: $\mat{P}$ for $n=2$ has $2!=2$ members (identity and the swap), $n=3$ has $3!=6$, and so on.
The product of two permutation matrices is another permutation matrix (do one shuffle, then the
other --- still a shuffle). The identity matrix is a permutation (the do-nothing shuffle). And every
permutation matrix has an inverse that is also a permutation matrix (reversing the shuffle undoes it).
A collection of operations with these properties --- closure under multiplication, associativity, an
identity element, and an inverse for every member --- is what mathematicians call a \emph{group}:
a self-contained algebraic system describing symmetries.

\begin{connectionbox}
The $n\times n$ permutation matrices form a group under multiplication, and it is an exact copy of
the \emph{symmetric group} $S_n$ --- the group of all rearrangements of $n$ symbols. This is the
historical birthplace of group theory (Galois, permuting the roots of a polynomial). So when you
multiply two permutation matrices you are composing two shuffles, and the noncommutativity you may
notice (doing swap-then-cycle differs from cycle-then-swap) is the same noncommutativity that makes
$S_3$ the smallest nonabelian group. In fact, by Cayley's theorem, every finite group can be
realized as a group of permutation matrices.
\end{connectionbox}

\section{The inverse of a permutation is its transpose}

Here the two halves of the chapter meet. To undo a shuffle you shuffle back, and the ``shuffle
back'' matrix is precisely the \emph{transpose} of the original.

\begin{theorem}[Permutations are orthogonal]\label{ch10:thm-perm-orth}
Every permutation matrix $\mat{P}$ satisfies
\[
  \mat{P}\T\mat{P} = \mat{I},
  \qquad\text{equivalently}\qquad
  \mat{P}\inv = \mat{P}\T.
\]
\end{theorem}

\begin{proof}
The columns of $\mat{P}$ are the standard basis vectors $\vec{e}_1,\dots,\vec{e}_n$ in \emph{some}
order; say column $j$ of $\mat{P}$ is $\vec{e}_{\sigma(j)}$ for some rearrangement $\sigma$. The
$(i,j)$ entry of $\mat{P}\T\mat{P}$ is the dot product of column $i$ with column $j$ of $\mat{P}$:
\[
  (\mat{P}\T\mat{P})_{ij} = \vec{e}_{\sigma(i)}\T\,\vec{e}_{\sigma(j)}.
\]
Standard basis vectors are orthonormal: $\vec{e}_a\T\vec{e}_b$ is $1$ when $a=b$ and $0$ otherwise.
Since $\sigma$ is a one-to-one rearrangement, $\sigma(i)=\sigma(j)$ exactly when $i=j$. So the dot
product is $1$ on the diagonal and $0$ off it. That is the identity: $\mat{P}\T\mat{P}=\mat{I}$.
Hence $\mat{P}\T$ is a left inverse of the square matrix $\mat{P}$, so it is the inverse.
\end{proof}

This is your \emph{first orthogonal matrix}. A square matrix $\mat{Q}$ is called \emph{orthogonal} when
$\mat{Q}\T\mat{Q}=\mat{I}$, which means its columns are \emph{orthonormal}: every column has length $1$
(normal) and every pair of distinct columns is perpendicular (orthogonal, with dot product zero).
Theorem~\ref{ch10:thm-perm-orth} says
permutation matrices are orthogonal, and they are the most transparent example: their columns are
literally the standard basis, shuffled, which is the most orthonormal a set of vectors can be.

\begin{example}[Transpose undoes the cycle]\label{ch10:ex-cycle-inverse}
Take the $3$-cycle from before and its transpose:
\[
  \mat{P}=\begin{bmatrix} 0 & 0 & 1 \\ 1 & 0 & 0 \\ 0 & 1 & 0 \end{bmatrix},
  \qquad
  \mat{P}\T=\begin{bmatrix} 0 & 1 & 0 \\ 0 & 0 & 1 \\ 1 & 0 & 0 \end{bmatrix}.
\]
Recall $\mat{P}$ sends $(x,y,z)\mapsto(z,x,y)$. Its transpose sends $(x,y,z)\mapsto(y,z,x)$ --- the
opposite cycle. Compose them:
\[
  \mat{P}\T\mat{P}
  = \begin{bmatrix} 0 & 1 & 0 \\ 0 & 0 & 1 \\ 1 & 0 & 0 \end{bmatrix}
    \begin{bmatrix} 0 & 0 & 1 \\ 1 & 0 & 0 \\ 0 & 1 & 0 \end{bmatrix}
  = \begin{bmatrix} 1 & 0 & 0 \\ 0 & 1 & 0 \\ 0 & 0 & 1 \end{bmatrix} = \mat{I}.
\]
Shuffle forward, shuffle back, home again. For a plain swap like
$\left[\begin{smallmatrix}0&1\\1&0\end{smallmatrix}\right]$ the matrix equals its own transpose, so
the swap is its own inverse --- which matches the picture: reflecting across $y=x$ twice does nothing.
\end{example}

\begin{aestheticbox}
Notice how the two halves of this chapter lock together. The transpose was defined so that
$\inner{\mat{Q}\vec{x}}{\mat{Q}\vec{y}} = \inner{\vec{x}}{\mat{Q}\T\mat{Q}\vec{y}}$.
Preserving every dot product therefore requires $\mat{Q}\T\mat{Q}=\mat{I}$, which immediately
forces $\mat{Q}\inv=\mat{Q}\T$. Permutations are the discrete prototype of this rule, preserving
lengths by shuffling axes; in Chapter~17, rotations and reflections will do the same thing
continuously.
\end{aestheticbox}

The geometric payoff is that orthogonal matrices --- and permutations in particular --- preserve
lengths. Here is the direct verification.

\begin{proposition}[Orthogonal matrices preserve length and angle]\label{ch10:prop-isometry}
If $\mat{Q}\T\mat{Q}=\mat{I}$, then for all $\vec{x},\vec{y}$,
\[
  (\mat{Q}\vec{x})\T(\mat{Q}\vec{y}) = \vec{x}\T\vec{y},
  \qquad\text{and in particular}\qquad
  \norm{\mat{Q}\vec{x}} = \norm{\vec{x}}.
\]
\end{proposition}

\begin{proof}
Use the adjoint relation (Theorem~\ref{ch10:thm-adjoint}) directly:
\[
  (\mat{Q}\vec{x})\T(\mat{Q}\vec{y})
  = \vec{x}\T\mat{Q}\T\mat{Q}\,\vec{y}
  = \vec{x}\T\mat{I}\,\vec{y}
  = \vec{x}\T\vec{y}.
\]
Setting $\vec{y}=\vec{x}$ gives $\norm{\mat{Q}\vec{x}}^2=\norm{\vec{x}}^2$, and lengths are
nonnegative, so $\norm{\mat{Q}\vec{x}}=\norm{\vec{x}}$.
\end{proof}

For a permutation this is obvious in hindsight: $\mat{P}\vec{x}$ has the same numbers as $\vec{x}$,
just in a different order, and reordering the entries cannot change $\sqrt{x_1^2+\cdots+x_n^2}$
(Figure~\ref{ch10:fig-isometry}). The
proposition says the abstract reason (orthogonality) agrees with the concrete one (same numbers,
reshuffled). The determinant of a permutation matrix is therefore $\pm 1$ --- it is $+1$ for an even
number of swaps and $-1$ for an odd number, the ``sign of the permutation'' that will reappear when
we build determinants.

\begin{figure}[H]
\centering
\begin{tikzpicture}[scale=1.0,every node/.style={font=\small}]
  % before
  \begin{scope}
    \draw[thin,gray!40] (-0.3,0)--(2.6,0);
    \draw[thin,gray!40] (0,-0.3)--(0,2.6);
    \draw[-{Latex[length=2.5mm]},vecblue,very thick] (0,0)--(2,1)
      node[midway,above left]{$\vec{x}$};
    \draw[densely dashed,gray] (2,0)--(2,1)--(0,1);
    \node[vecblue,right] at (2.05,1.05) {$(2,1)$};
    \node[gray!60] at (1.15,-0.55) {$\norm{\vec{x}}=\sqrt5$};
  \end{scope}
  \draw[-{Latex[length=3mm]},thick] (3.0,1)--(4.3,1)
    node[midway,above]{$\mat{P}=\left[\begin{smallmatrix}0&1\\1&0\end{smallmatrix}\right]$};
  % after: swapped to (1,2)
  \begin{scope}[shift={(5.0,0)}]
    \draw[thin,gray!40] (-0.3,0)--(2.6,0);
    \draw[thin,gray!40] (0,-0.3)--(0,2.6);
    \draw[densely dashed,trapred!60] (-0.2,-0.2)--(2.4,2.4);
    \draw[-{Latex[length=2.5mm]},vecred,very thick] (0,0)--(1,2)
      node[midway,above left]{$\mat{P}\vec{x}$};
    \draw[densely dashed,gray] (1,0)--(1,2)--(0,2);
    \node[vecred,above] at (1.05,2.05) {$(1,2)$};
    \node[gray!60] at (1.15,-0.55) {$\norm{\mat{P}\vec{x}}=\sqrt5$};
  \end{scope}
\end{tikzpicture}
\caption{A permutation only relabels coordinates, so it lands $\vec{x}=(2,1)$ on
$\mat{P}\vec{x}=(1,2)$, the reflection across $y=x$. Same two numbers, same length
$\sqrt{2^2+1^2}=\sqrt5$. Orthogonality ($\mat{P}\T\mat{P}=\mat{I}$) is the abstract name for this
length preservation.}
\label{ch10:fig-isometry}
\end{figure}

\begin{walkbox}
Take a walk and turn this over. A permutation matrix preserves all lengths and angles, so as a
rigid motion it is some rotation or reflection of space. Question: which permutations are pure
rotations (determinant $+1$) and which include a reflection (determinant $-1$)? Try the $2\times 2$
swap (is exchanging $x$ and $y$ a rotation, or a flip?) and the $3$-cycle. Can a single transposition
ever be a rotation of the plane? What does its determinant tell you?
\end{walkbox}

\begin{recoverybox}
\textbf{Pause and think:} Why must every orthogonal matrix have determinant equal to $+1$ or $-1$?
\emph{Answer:} Take the determinant of both sides of $\mat{Q}\T\mat{Q}=\mat{I}$.
Using $\det(\mat{A}\mat{B})=\det(\mat{A})\det(\mat{B})$ and $\det(\mat{Q}\T)=\det(\mat{Q})$,
we find $\big(\det(\mat{Q})\big)^2 = \det(\mat{I}) = 1$, which forces $\det(\mat{Q}) = \pm 1$.
A determinant of $+1$ means the transformation preserves orientation (like a pure rotation);
a determinant of $-1$ means orientation is reversed (including a reflection).
\end{recoverybox}

\section{Putting transposes and permutations together}

The two ideas interlock in elimination, which plays a central role in subsequent chapters.
When Gaussian elimination needs to avoid a zero pivot, it swaps rows --- a left multiplication
by a permutation $\mat{P}$. The clean statement of elimination ``with row exchanges'' is the
factorization
\[
  \mat{P}\mat{A} = \mat{L}\mat{U},
\]
where $\mat{P}$ is a permutation collecting all the swaps, $\mat{L}$ is lower triangular, and
$\mat{U}$ is upper triangular. Because $\mat{P}\inv=\mat{P}\T$, you can move the permutation to the
other side for free: $\mat{A}=\mat{P}\T\mat{L}\mat{U}$. No inverse needs computing; the transpose
\emph{is} the inverse. This is the practical payoff of Theorem~\ref{ch10:thm-perm-orth} and it is why
numerical libraries store a permutation as a list of swaps rather than as a matrix --- transposing a
list of swaps is just reading it backward.

\begin{example}[A row swap rescuing a pivot]\label{ch10:ex-pivot}
Consider
\[
  \mat{A}=\begin{bmatrix} 0 & 2 \\ 3 & 4 \end{bmatrix}.
\]
The top-left pivot is $0$, so elimination cannot proceed. Swap the rows with
$\mat{P}=\left[\begin{smallmatrix}0&1\\1&0\end{smallmatrix}\right]$:
\[
  \mat{P}\mat{A}
  = \begin{bmatrix} 0&1\\1&0 \end{bmatrix}\begin{bmatrix} 0 & 2 \\ 3 & 4 \end{bmatrix}
  = \begin{bmatrix} 3 & 4 \\ 0 & 2 \end{bmatrix},
\]
which now has a healthy pivot $3$ in the top-left and is already upper triangular (here $\mat{L}=\mat{I}$,
$\mat{U}=\mat{P}\mat{A}$). To recover $\mat{A}$ itself, multiply by $\mat{P}\T=\mat{P}$ on the left
again --- the swap is its own undo (Figure~\ref{ch10:fig-pivot-rescue}).
\end{example}

\begin{figure}[H]
\centering
\begin{tikzpicture}[scale=0.9,every node/.style={font=\small}]
  % Left matrix A
  \begin{scope}[shift={(0,0)}]
    % Matrix brackets
    \draw[thick,gray!60] (0.0,1.05) -- (-0.25,1.05) -- (-0.25,-1.05) -- (0.0,-1.05);
    \draw[thick,gray!60] (2.0,1.05) -- (2.25,1.05) -- (2.25,-1.05) -- (2.0,-1.05);
    % Row 1
    \draw[trapred,thick,rounded corners,fill=trapred!8] (-0.1,0.15) rectangle (2.1,0.95);
    \node at (0.4,0.55) {$\mathbf{0}$};
    \node at (1.5,0.55) {$2$};
    \node[trapred,above,font=\footnotesize] at (1.0,1.05) {row 1 (pivot $0$: blocked)};
    % Row 2
    \draw[targetgreen,thick,rounded corners,fill=targetgreen!8] (-0.1,-0.95) rectangle (2.1,-0.15);
    \node at (0.4,-0.55) {$\mathbf{3}$};
    \node at (1.5,-0.55) {$4$};
    \node[targetgreen,below,font=\footnotesize] at (1.0,-1.05) {row 2 (pivot $3$: ready)};
    \node[below=1.4cm] at (1.0,-0.15) {original $\mat{A}$};
  \end{scope}

  % Middle arrow
  \draw[-{Latex[length=3mm]},thick,sourceblue] (2.7,0) -- (5.9,0)
    node[midway,above,font=\small] {$\mat{P}$ (swap $R_1 \leftrightarrow R_2$)};

  % Right matrix PA
  \begin{scope}[shift={(6.6,0)}]
    % Matrix brackets
    \draw[thick,gray!60] (0.0,1.05) -- (-0.25,1.05) -- (-0.25,-1.05) -- (0.0,-1.05);
    \draw[thick,gray!60] (2.0,1.05) -- (2.25,1.05) -- (2.25,-1.05) -- (2.0,-1.05);
    % Row 1 (former Row 2)
    \draw[targetgreen,thick,rounded corners,fill=targetgreen!8] (-0.1,0.15) rectangle (2.1,0.95);
    \node at (0.4,0.55) {$\mathbf{3}$};
    \node at (1.5,0.55) {$4$};
    \node[targetgreen,above,font=\footnotesize] at (1.0,1.05) {row 1 (active pivot $3$)};
    % Row 2 (former Row 1)
    \draw[trapred,thick,rounded corners,fill=trapred!8] (-0.1,-0.95) rectangle (2.1,-0.15);
    \node at (0.4,-0.55) {$0$};
    \node at (1.5,-0.55) {$\mathbf{2}$};
    \node[trapred,below,font=\footnotesize] at (1.0,-1.05) {row 2 (active pivot $2$)};
    \node[below=1.4cm] at (1.0,-0.15) {$\mat{P}\mat{A} = \mat{U}$};
  \end{scope}
\end{tikzpicture}
\caption{Permutation matrix as an elimination pivot rescue. In $\mat{A}$ (left), the top-left pivot
is $0$, so Gaussian elimination cannot proceed without division by zero. Multiplying on the left by
the row-swap permutation $\mat{P}$ exchanges row 1 and row 2 (right). The resulting matrix $\mat{P}\mat{A}$
has a valid pivot $3$ in place and is already in upper-triangular echelon form $\mat{U}$.}
\label{ch10:fig-pivot-rescue}
\end{figure}

\begin{trapbox}
A subtle but common error: people write $\mat{A}=\mat{L}\mat{U}\mat{P}$ or sprinkle the permutation
on the right. The swaps in elimination act on \emph{rows}, so $\mat{P}$ must sit on the \emph{left}:
$\mat{P}\mat{A}=\mat{L}\mat{U}$. Putting it on the right would permute columns, which is a different
operation (it reorders the unknowns, not the equations). Left for rows, right for columns --- keep
that straight and the factorization comes out right.
\end{trapbox}

\begin{formalbox}
Collecting the main identities of this chapter, all for real matrices:
\begin{itemize}
  \item $(\mat{A}\mat{B})\T=\mat{B}\T\mat{A}\T$ and $(\mat{A}\T)\T=\mat{A}$.
  \item $\inner{\mat{A}\vec{x}}{\vec{y}}=\inner{\vec{x}}{\mat{A}\T\vec{y}}$ for all $\vec{x},\vec{y}$
        (the adjoint relation, which characterizes $\mat{A}\T$ uniquely).
  \item $\mat{A}\T\mat{A}$ and $\mat{A}\mat{A}\T$ are symmetric; the $(i,j)$ entry of $\mat{A}\T\mat{A}$
        is $\vec{a}_i\T\vec{a}_j$ (the Gram matrix of the columns).
  \item A permutation matrix has one $1$ per row and column; it is orthogonal,
        $\mat{P}\T\mat{P}=\mat{I}$, so $\mat{P}\inv=\mat{P}\T$.
  \item Orthogonal matrices ($\mat{Q}\T\mat{Q}=\mat{I}$) preserve dot products, lengths, and angles;
        their determinant is $\pm 1$.
\end{itemize}
\end{formalbox}

\subsection{Looking forward: from matrix arithmetic to vector spaces}

Throughout the first ten chapters, our calculations have revolved around concrete systems of equations
and grids of numbers: row reductions, matrix products, inverses, factorizations like $\mat{L}\mat{U}$
and $\mat{P}\mat{A}=\mat{L}\mat{U}$, and now the transpose $\mat{A}\T$.

Yet every rectangular matrix $\mat{A}$ of shape $m\times n$ connects two geometric worlds at once:
it pushes vectors forward from $\R^n$ to $\R^m$, while its transpose $\mat{A}\T$ pulls vectors backward
from $\R^m$ to $\R^n$. Some inputs get stretched into new directions; others get squashed completely to zero.
What are the shapes of the regions formed by all reachable outputs? What are the shapes of the directions
that get crushed?

To answer these questions, we must leave behind mechanical entry-by-entry arithmetic and step into the
conceptual foundation of the entire subject: \emph{vector spaces and subspaces}. In Chapter~11, we begin
Part~2 by exploring these geometric worlds where vectors live, add, and scale. That path leads directly to
Gilbert Strang's celebrated ``Big Picture of Linear Algebra'' in Chapter~15, where the transpose $\mat{A}\T$
returns to claim center stage: proving that the directions $\mat{A}$ squashes to zero are strictly perpendicular
to the directions $\mat{A}\T$ can reach.

\begin{historybox}
The transpose began as a humble bookkeeping device for organizing algebraic formulas. In 1855, the British
mathematician Arthur Cayley introduced matrices to study how algebraic equations transform when you change
geometric coordinates. Cayley was tired of writing out endless systems of equations variable by variable;
he realized that flipping a rectangular table across its main diagonal turned row equations into column
equations, calling the result the ``transposed matrix.'' For decades, textbooks treated Cayley's transpose
as little more than a convenient typing shortcut.

The true depth of the transpose emerged in the early twentieth century, when David Hilbert, Stefan Banach,
and Frigyes Riesz developed \emph{functional analysis} --- the study of infinite-dimensional spaces where
the ``vectors'' are continuous functions and the ``matrices'' are differential equations or integrals.
In continuous spaces, there are no grid tiles to flip across a diagonal. How do you transpose a
derivative? Riesz and Banach answered by looking at \emph{dual spaces} --- collections of linear measurement
probes that turn vectors into scalars, like measuring the temperature or velocity of a system at a point.
The \emph{Riesz Representation Theorem} showed that in any space equipped with an inner product (a dot product
for functions), every linear measurement probe corresponds to taking the inner product with a unique vector.
This led to the modern definition of the \emph{adjoint}: the transpose $\mat{A}\T$ is not defined by its
table of entries, but by the fundamental balance $\inner{\mat{A}\vec{x}}{\vec{y}} = \inner{\vec{x}}{\mat{A}\T\vec{y}}$.
The finite matrix transpose is simply the Euclidean coordinate shadow of this universal geometric relationship.

In numerical computation, permutation matrices sparked a fierce debate during the dawn of electronic computers
in the 1950s. Pioneers like Alan Turing and John von Neumann feared that Gaussian elimination was fundamentally
unstable, warning that tiny floating-point rounding errors would cascade and ruin the solution. In 1961, the
British numerical analyst James Wilkinson proved that keeping pivots large via row swaps --- Gaussian elimination
with partial pivoting, $\mat{P}\mat{A} = \mat{L}\mat{U}$ --- controls this error growth. Yet Wilkinson uncovered
a famous controversy: in the theoretical worst case, partial pivoting can amplify errors by a factor of
$2^{n-1}$, an exponential worst-case bound. To this day, mathematicians and computer scientists analyze why
this worst-case behavior almost never occurs in practice: on real-world engineering problems, $\mat{P}\mat{A}=\mat{L}\mat{U}$
behaves with reliable numerical stability.

Through the transpose we define symmetric matrices ($\mat{A}\T = \mat{A}$), which represent observables in quantum
mechanics, stress tensors in continuum mechanics, and covariance matrices in statistics, as well as the
Gram matrix $\mat{A}\T\mat{A}$, which is always symmetric and \emph{positive semidefinite} --- meaning it
never assigns negative energy to any vector, $\vec{x}\T(\mat{A}\T\mat{A})\vec{x}=\norm{\mat{A}\vec{x}}^2 \ge 0$.
This identity is the foundation of least-squares data fitting, the normal equations, and backpropagation
in deep learning.

Yet deep questions about permutations and transposes remain active research frontiers today. The celebrated
\emph{Graph Isomorphism} problem asks whether two network graphs share the exact same connection structure ---
which is mathematically identical to asking whether their connection matrices $\mat{A}$ and $\mat{B}$ are
related by a permutation matrix: $\mat{P}\T\mat{A}\mat{P} = \mat{B}$. In 2015, László Babai shook theoretical
computer science by proving that graph isomorphism can be solved in quasipolynomial time, but whether a truly
efficient polynomial-time algorithm exists remains a famous open question. Meanwhile, in massive scientific
simulations, finding the ideal row permutation $\mat{P}$ to minimize ``fill-in'' (preventing zero entries from
accidentally turning into expensive nonzeros during $\mat{P}\mat{A}=\mat{L}\mat{U}$) is known to be NP-complete,
driving ongoing research into heuristic graph algorithms for supercomputers.
\end{historybox}

\section{Exercises}

\begin{exercise}[Warmup flips]\label{ch10:ex-warmup}
Write down the transpose of each matrix and state its shape.
\[
  \mat{A}=\begin{bmatrix} 2 & -1 & 0 \\ 5 & 3 & 7 \end{bmatrix},\qquad
  \mat{B}=\begin{bmatrix} 1 & 4 \\ 4 & 9 \end{bmatrix},\qquad
  \vec{c}=\colvec{3\\-2\\6}.
\]
Which of these (if any) is symmetric, and how can you tell at a glance?
\end{exercise}

\begin{exercise}[Reversal by hand]\label{ch10:ex-reversal}
Let $\mat{A}=\left[\begin{smallmatrix}2&0\\1&3\end{smallmatrix}\right]$ and
$\mat{B}=\left[\begin{smallmatrix}1&5\\0&2\end{smallmatrix}\right]$. Compute $\mat{A}\mat{B}$,
$(\mat{A}\mat{B})\T$, $\mat{B}\T\mat{A}\T$, and $\mat{A}\T\mat{B}\T$. Confirm which two agree and
explain in one sentence why the fourth one is the odd one out.
\end{exercise}

\begin{exercise}[Dot products as transposes]\label{ch10:ex-dotprod}
Let $\vec{u}=\colvec{1\\2\\-2}$ and $\vec{v}=\colvec{3\\0\\1}$. Compute $\vec{u}\T\vec{v}$,
$\vec{v}\T\vec{u}$, $\norm{\vec{u}}^2=\vec{u}\T\vec{u}$, and the outer product $\vec{u}\vec{v}\T$
(a $3\times 3$ matrix). Is the outer product symmetric? Why or why not?
\end{exercise}

\begin{exercise}[Building Gram matrices]\label{ch10:ex-gram}
For $\mat{A}=\left[\begin{smallmatrix}1&2\\0&1\\1&0\end{smallmatrix}\right]$ compute $\mat{A}\T\mat{A}$
and $\mat{A}\mat{A}\T$. Verify both are symmetric. Identify which entry of $\mat{A}\T\mat{A}$ is the
dot product of the two columns of $\mat{A}$, and check it by computing that dot product directly.
\end{exercise}

\begin{exercise}[Symmetric plus antisymmetric]\label{ch10:ex-split}
Split $\mat{A}=\left[\begin{smallmatrix}4&1\\7&2\end{smallmatrix}\right]$ into its symmetric part
$\tfrac12(\mat{A}+\mat{A}\T)$ and its antisymmetric part $\tfrac12(\mat{A}-\mat{A}\T)$. Add them back
to confirm you recover $\mat{A}$. What must the diagonal of any antisymmetric matrix be, and why?
\end{exercise}

\begin{exercise}[Permutations and their actions]\label{ch10:ex-permaction}
Let $\mat{P}=\left[\begin{smallmatrix}0&0&1\\0&1&0\\1&0&0\end{smallmatrix}\right]$. Describe in words
what $\mat{P}\vec{x}$ does to $\vec{x}=(x_1,x_2,x_3)$. Compute $\mat{P}\T$ and verify
$\mat{P}\T\mat{P}=\mat{I}$. Is this $\mat{P}$ its own inverse? Explain using the cycle structure.
\end{exercise}

\begin{exercise}[Left versus right]\label{ch10:ex-leftright}
With $\mat{P}=\left[\begin{smallmatrix}0&1\\1&0\end{smallmatrix}\right]$ and an arbitrary
$\mat{M}=\left[\begin{smallmatrix}a&b\\c&d\end{smallmatrix}\right]$, compute $\mat{P}\mat{M}$ and
$\mat{M}\mat{P}$. State precisely what each does to $\mat{M}$ and connect it to row swaps versus
column swaps in elimination.
\end{exercise}

\begin{exercise}[Adjoint by example]\label{ch10:ex-adjoint}
Let $\mat{A}=\left[\begin{smallmatrix}1&2\\0&1\\3&0\end{smallmatrix}\right]$,
$\vec{x}=\colvec{1\\-1}$, and $\vec{y}=\colvec{2\\0\\1}$. Compute $\mat{A}\vec{x}$ and
$\inner{\mat{A}\vec{x}}{\vec{y}}$; then compute $\mat{A}\T\vec{y}$ and $\inner{\vec{x}}{\mat{A}\T\vec{y}}$.
Confirm the two numbers agree, illustrating Theorem~\ref{ch10:thm-adjoint}.
\end{exercise}

\begin{exercise}[Counting and composing]\label{ch10:ex-count}
How many $3\times 3$ permutation matrices are there? Pick the swap of rows $1$ and $2$, call it
$\mat{P}_{12}$, and the swap of rows $2$ and $3$, call it $\mat{P}_{23}$. Compute $\mat{P}_{12}\mat{P}_{23}$
and $\mat{P}_{23}\mat{P}_{12}$. Are they equal? What does your answer say about whether the group of
permutation matrices is commutative?
\end{exercise}

\begin{exercise}[Orthogonality preserves geometry]\label{ch10:ex-isometry}
Let $\mat{P}$ be any $3\times 3$ permutation matrix and $\vec{x}=(1,-2,2)$. Without choosing a
specific $\mat{P}$, explain why $\norm{\mat{P}\vec{x}}=3$ no matter which permutation you pick. Then
prove in general that if $\mat{Q}\T\mat{Q}=\mat{I}$ and $\mat{R}\T\mat{R}=\mat{I}$ then the product
$\mat{Q}\mat{R}$ is also orthogonal.
\end{exercise}

\begin{exercise}[A proof to finish]\label{ch10:ex-proof}
Prove that if $\mat{A}$ is any $m\times n$ matrix and $\mat{A}\T\mat{A}=\mat{0}$ (the zero matrix),
then $\mat{A}=\mat{0}$. \emph{Hint:} look at the diagonal entries of $\mat{A}\T\mat{A}$ and recall
what they equal in terms of the columns of $\mat{A}$. In particular, $\mat{A}\T\mat{A}$
cannot collapse to zero unless $\mat{A}$ already was zero, which is the seed of why
$\mat{A}\T\mat{A}$ is invertible exactly when $\mat{A}$ has independent columns.
\end{exercise}

\begin{exercise}[Elimination with row swaps: \texorpdfstring{$\mat{P}\mat{A}=\mat{L}\mat{U}$}{PA=LU}]\label{ch10:ex-pa-lu}
Consider the matrix
\[
  \mat{A} = \begin{bmatrix} 0 & 1 & 3 \\ 2 & 1 & 1 \\ 4 & 5 & 2 \end{bmatrix}.
\]
\begin{enumerate}[label=(\alph*)]
  \item Explain why standard Gaussian elimination cannot begin without swapping rows.
  \item Let $\mat{P}$ be the permutation matrix that exchanges row 1 and row 2. Write down $\mat{P}$
        and compute the product $\mat{P}\mat{A}$.
  \item Perform elimination on $\mat{P}\mat{A}$ to find a unit lower-triangular matrix $\mat{L}$ and
        an upper-triangular matrix $\mat{U}$ such that $\mat{P}\mat{A} = \mat{L}\mat{U}$.
  \item Using $\mat{P}\inv=\mat{P}\T$, write down the factorization of the original matrix as
        $\mat{A} = \mat{P}\T\mat{L}\mat{U}$, and verify by multiplication that it recovers $\mat{A}$.
\end{enumerate}
\end{exercise}

\begin{exercise}[Permutation similarity and graph isomorphism]\label{ch10:ex-graph-iso}
Let a $3$-node network (graph) have connections given by the symmetric adjacency matrix
\[
  \mat{A} = \begin{bmatrix} 0 & 1 & 0 \\ 1 & 0 & 1 \\ 0 & 1 & 0 \end{bmatrix},
\]
where entry $a_{ij}=1$ indicates an undirected edge connecting node $i$ and node $j$
(a line network: node 1 is connected to node 2, which is connected to node 3).
\begin{enumerate}[label=(\alph*)]
  \item Let $\mat{P}=\left[\begin{smallmatrix}0&1&0\\1&0&0\\0&0&1\end{smallmatrix}\right]$ be the
        permutation that exchanges node 1 and node 2. Compute $\mat{B} = \mat{P}\T\mat{A}\mat{P}$.
  \item Explain why left-multiplication by $\mat{P}\T$ permutes the rows while right-multiplication
        by $\mat{P}$ permutes the columns, so that $\mat{P}\T\mat{A}\mat{P}$ renames both source and
        target of every edge consistently.
  \item Verify that $\mat{B}$ is also symmetric. Which node is now the central hub connected to the
        other two? Conclude that $\mat{B}$ represents the exact same network structure, merely with
        the node labels renumbered.
\end{enumerate}
\end{exercise}
